% Candidate section; requires amsmath, amssymb, amsthm.
% Status: computer-assisted proof candidate, not peer reviewed.
% Do not claim human hand-verification of every arithmetic step.
\section{Asymptotics of the two-colour threshold}\label{sec:K-asymptotic}
\noindent\emph{Status and provenance.}
The results below are a computer-assisted proof candidate, not peer reviewed.
AI systems helped design, calculate, independently check, and write the proofs;
the human author retains responsibility for any eventual submission.
In particular, the upper bound invokes the balanced-window and parity-window
lemmas, whose proof and independent rechecking are described separately.

For $k\ge\ell\ge6$, let
\[
 F(k,\ell)=k^2+\left(\binom\ell2+1\right)k+\ell(\ell-2),
 \qquad a_* = \frac{\sqrt2-1}{2},
\]
and let $K(\ell)$ denote the least integer $K$ such that
$S_2(k;\ell)=F(k,\ell)$ for every $k\ge K$.
We use $e=k-\ell$.

\begin{theorem}[Asymptotic threshold; proof candidate]\label{thm:K-asymptotic}
For every integer $\ell\ge50$,
\[
 (1+a_*)\ell-4 < K(\ell) < (1+a_*)\ell+5.
\]
For $6\le\ell\le49$, the finite certificate table gives
$K(\ell)<(1+a_*)\ell+9$.
Consequently
\[
 K(\ell)=\frac{1+\sqrt2}{2}\,\ell+O(1).
\]
The upper bound uses the separately established balanced-window and
parity-window lemmas for all $\ell\ge6$, $e\ge1$.
\end{theorem}

\begin{proof}[Lower bound]
For $1\le j\le\ell-1$ set
\[
 m=\frac{\ell(\ell+1)}2+\ell j+e(j+1)-1,
 \qquad q=\ell-j,
\]
\[
 V_3=\frac{j(j+1)}2+q(m+1)+\binom q2+e.
\]
Colour $[1,V_3-1]$ red on $[1,j]\cup[m+1,V_3-1]$ and blue on
$[j+1,m]$. Any blue $k$-sum using exactly $\ell$ distinct summand
values is at least
$\sum_{i=1}^\ell(j+i)+e(j+1)=m+1$, hence not blue.
A red $k$-sum has at most $j$ distinct small summands and at least
$q$ distinct summands at least $m+1$; it therefore is at least
$\sum_{i=1}^{j}i+\sum_{i=1}^{q}(m+i)+e=V_3$ and hence
cannot be red. The colouring is good, so $S_2(k;\ell)\ge V_3$.

Set $G(j)=V_3-F(k,\ell)=-Aj^2+\beta j+\gamma$,
where $A=e+\ell-1>0$. Its vertex is
\[
 j_*=\frac{(\ell-1)(2e+\ell-2)}{4(e+\ell-1)}.
\]
For $e_1=a_*\ell-4$ and all $\ell\ge50$ it lies within
$[1,\ell-1]$. Rounding the vertex to an integer loses at most
$A/4$, so it suffices to check
\[
 4A\gamma+\beta^2-A^2-4A\ge0.
\]
An exact symbolic expansion in $u=\ell-50\ge0$, $e=e_1$ has
nonnegative coefficients (the coefficient list and checker are
included with the archived independent audit). Moreover
$\partial G/\partial e<0$ for all admissible $j$ and $e\ge1$.
Thus whenever $1\le e\le a_*\ell-4$, some integer $j$ satisfies
$V_3>F(\ell+e,\ell)$, excluding that $k$ from the eventual equality
range. Hence $K(\ell)>(1+a_*)\ell-4$.
\end{proof}

\begin{lemma}[T1-fill, sparse windows]\label{lem:T1-fill}
Suppose the initial interval $[1,2\ell]$ has a colour class $D$
with $2\le j=|D|\le\ell-2$ and opposite class $M$.
Put $q=\ell-j$ and let $S$ be the least $\ell$ members of $M$.
There exist $y,r',d\in S$ and $b\in D$ with $b=y+r'-d$.
For every $e\ge1$, the six families of forced values $L_1,L_2,Z_i,M_v,w,w'$ defined below constitute
an impossible colouring whenever $w\le F(\ell+e,\ell)$.
\end{lemma}
\begin{proof}
If no triple existed, $S$ would be closed under $x+y-z$ whenever
the result lay in $[1,\max S]$. Write $m_1<m_2$ for the first
members and $\delta=m_2-m_1$. A standard up-and-down translation
argument yields $S=(m_1+\delta\mathbb Z)\cap[1,\max S]$,
with $\delta\ge m_1$. For $\delta=1$, the minority colour has no
value below $\max S$, contrary to the assumed nonempty intersection.
For $\delta\ge2$,
$\max S\ge m_1+2(\ell-1)>\ell+j\ge\max S$, a contradiction.
If the minority class starts above $\max S$, then $S=[1,\ell]$,
and choosing $y=\ell$, $r'=b+1-\ell$, $d=1$ works for any
minority $b\le2\ell-1$, which exists since $j\ge2$.

Set $R'=S\setminus\{y\}$, $u=\min R'$,
$r=\min(S\setminus\{r'\})$,
$z_1,\ldots,z_{q-1}$ the $q-1$ least members of
$S\setminus\{r,r'\}$, $d_b=\min(D\setminus\{b\})$,
and $\Sigma_0=\sum S+(e-1)u$. All of
\[
 L_1=\Sigma_0+r,\quad L_2=\Sigma_0+r',\quad
 Z_i=\Sigma_0+z_i
\]
are $M$-colour $k$-sums and therefore are forced to $D$.
Writing $B_3=(D\setminus\{b\})\cup\{Z_i:1\le i<q\}$,
we have $|B_3|=\ell-2$, and
\[
 M_v=\sum B_3+b+L_1+e\,d_b,
\quad w=\sum R'+M_v+(e-1)u+d,
\]
\[
 w'=\sum B_3+L_1+L_2+e\,d_b.
\]
The first sum forces $M_v$ to $M$, the second forces $w$ to $D$,
and the third forces $w'$ to $M$. All input sets have precisely
$\ell$ distinct members, each output exceeds its summands, and
$w-w'=b+d-y-r'=0$. In this chain $w$ is the largest integer;
if $w\le F$, a contradiction follows.
\end{proof}

\begin{proof}[Upper bound in Theorem~\ref{thm:K-asymptotic}]
Assume a good colouring of $[1,F(\ell+e,\ell)]$.
The balanced-window lemma excludes the case in which both colours
occur at least $\ell-1$ times in $[1,2\ell]$, including the parity
window (proved by an odd/even case split). If the window is all one
colour, the consecutive-interval $k$-sum lemma applied twice gives
a contradiction for all $e\ge1$. If one colour occurs only once,
a three-stage consecutive-sum argument excludes the window for
$\ell\ge7$, $e\ge3$ (or $\ell=6,e\ge4$).

It remains to take $2\le j\le\ell-2$ and apply
Lemma~\ref{lem:T1-fill}. Bounding its final value $w$ above using
the least $\ell$ majority points and the greatest possible positions
of $j$ minority points gives two explicit quadratic polynomials
$P_{\rm I}(j)=F-w_{\rm bound}^{\rm I}$ and
$P_{\rm II}(j)=F-w_{\rm bound}^{\rm II}$, with leading
coefficients $\ell+1$ and $\ell+e$, respectively. The exact
arithmetic checks attached to this preprint prove
\[
 4\alpha\gamma-\beta^2\ge0
 \quad\text{at }e=a_*\ell+4,\quad\ell\ge50
\]
for both polynomials $P(j)=\alpha j^2+\beta j+\gamma$.
The verification is by expansion in $u=\ell-50$ with nonnegative
coefficients in $\mathbb Q(\sqrt2)$; $P$ is increasing in $e$
throughout the required range. Therefore every integer
$e\ge\lceil a_*\ell+4\rceil$ rules out all sparse windows,
so $K(\ell)\le\ell+\lceil a_*\ell+4\rceil
<(1+a_*)\ell+5$.
The finite interval $6\le\ell\le49$ is checked by exhaustive
integer evaluations of the same upper-bound polynomials; its largest
additive overhead is $8.5502526\ldots<9$.
\end{proof}

\paragraph{Relaxed variant.}
The same asymptotic conclusion is conditional on interpreting the
relaxed condition as at least $\ell$ distinct \emph{summand} values,
i.e., $c(\mathcal S)\ge\ell+1$ including the output.
Gaiser's introduction uses this interpretation but its Section~6
contains $c(\mathcal S)\ge\ell$; the discrepancy requires
clarification before publication.

\paragraph{Verification statement.}
The independent review used newly written exact-integer and
$\mathbb Q(\sqrt2)$ symbolic programs, selected exhaustive finite
window scans and explicit sum certificates. This is not a formally
verified theorem in a proof assistant and has not been peer reviewed.
